% ---------------------------------------------------------------------------
% listings.tex -- source-code and terminal-session typesetting.
% Same palette and mechanism as the SystemVerilog cookbook, trimmed to the
% languages this booklet actually needs: shell sessions, Makefiles, plain
% output, and the .core/.hjson descriptor files.
% ---------------------------------------------------------------------------

% --- the palette (identical to the SystemVerilog cookbook) -----------------
\definecolor{codebg}      {HTML}{FAFAF7}
\definecolor{codeframe}   {HTML}{D9D6CC}
\definecolor{codekey}     {HTML}{0B5394}   % keywords / headings
\definecolor{codekeyb}    {HTML}{7B1FA2}   % secondary keywords
\definecolor{codecomment} {HTML}{5C7A5C}
\definecolor{codestring}  {HTML}{A0522D}
\definecolor{codenumber}  {HTML}{9A9A93}   % line numbers, muted text
\definecolor{accent}      {HTML}{1B5E20}
\definecolor{accentalt}   {HTML}{8B2500}
\definecolor{warnbg}      {HTML}{FDF3F2}
\definecolor{warnframe}   {HTML}{C0392B}
\definecolor{tipbg}       {HTML}{F1F8F1}
\definecolor{tipframe}    {HTML}{2E7D32}
\definecolor{notebg}      {HTML}{F0F4FA}
\definecolor{noteframe}   {HTML}{1F4E79}

% ---------------------------------------------------------------------------
% Base style shared by every language.
% ---------------------------------------------------------------------------
\lstdefinestyle{cookbookbase}{%
  basicstyle      = \ttfamily\footnotesize,
  keywordstyle    = \color{codekey}\bfseries,
  keywordstyle    = [2]\color{codekeyb}\bfseries,
  commentstyle    = \color{codecomment}\itshape,
  stringstyle     = \color{codestring},
  numberstyle     = \tiny\color{codenumber},
  numbers         = none,
  showstringspaces= false,
  showspaces      = false,
  showtabs        = false,
  tabsize         = 2,
  breaklines      = true,
  breakatwhitespace = false,
  breakindent     = 1em,
  postbreak       = \mbox{\textcolor{codenumber}{$\hookrightarrow$}\space},
  columns         = fullflexible,
  keepspaces      = true,
  upquote         = true,
  frame           = leftline,
  framerule       = 1.1pt,
  rulecolor       = \color{codeframe},
  framexleftmargin= 3pt,
  xleftmargin     = 14pt,
  aboveskip       = 1.1\medskipamount,
  belowskip       = 1.1\medskipamount,
  captionpos      = b,
  escapeinside    = {(*@}{@*)},
  moredelim       = **[is][\color{accentalt}\bfseries]{@@}{@@},
  backgroundcolor = \color{codebg},
}

% Shell / make command-line sessions. `$' marks a prompt; comments start `#'.
\lstdefinestyle{shstyle}{%
  style=cookbookbase,
  language=bash,
  morekeywords={fusesoc,make},
  keywordstyle=\color{codekey}\bfseries,
}

% Plain terminal output, file contents, logs.
\lstdefinestyle{plainstyle}{%
  style=cookbookbase,
  language={},
  keywordstyle=\color{black},
}

% Makefiles: targets, variables, the handful of GNU make directives used here.
\lstdefinestyle{makestyle}{%
  style=cookbookbase,
  language=make,
  backgroundcolor=\color{codebg},
}

% .core (FuseSoC, YAML-ish) and .hjson (vendoring descriptors) files. Neither
% language is built into `listings'; a light generic style with `#' and `//'
% comments recognised, plus the field names this booklet actually uses
% highlighted as keywords, is enough to read at booklet length.
\lstdefinestyle{corestyle}{%
  style=cookbookbase,
  language={},
  keywordstyle=\color{codekey}\bfseries,
  morekeywords={name,description,filesets,depend,targets,tools,files,
                file_type,toplevel,default_tool,default,rtl,tb,
                systemVerilogSource,upstream,url,rev,target_dir,
                exclude_from_upstream,verilator_options,vlog_options,
                mode,script_dir,dc_script,report_dir,libs},
  morecomment=[l]{//},
  morecomment=[l]{\#},
}

% SystemVerilog and C, used only where this booklet quotes real RTL/header
% excerpts (the register\_interface / regtool chapter) -- built-in `listings'
% languages, tuned with the handful of SV-only keywords they lack. Nothing
% like the full dialect/UVM highlighting of the SystemVerilog cookbook: this
% booklet only ever shows short, already-generated or already-working
% excerpts, never live coding.
\lstdefinestyle{svstyle}{%
  style=cookbookbase,
  language=Verilog,
  morekeywords={logic,bit,byte,shortint,int,longint,string,typedef,struct,
                packed,always_comb,always_ff,import,package,endpackage,
                interface,modport,parameter,localparam,automatic,unique,
                priority,return,break,continue,foreach},
  keywordstyle=[2]\color{codekeyb}\bfseries,
}

\lstdefinestyle{cppstyle}{%
  style=cookbookbase,
  language=C++,
  morekeywords={uint8_t,uint16_t,uint32_t,uint64_t,int8_t,int16_t,int32_t,
                int64_t,size_t},
}

\lstnewenvironment{shcode}  [1][]{\lstset{style=shstyle,   #1}}{}
\lstnewenvironment{txtcode} [1][]{\lstset{style=plainstyle,#1}}{}
\lstnewenvironment{makecode}[1][]{\lstset{style=makestyle, #1}}{}
\lstnewenvironment{corecode}[1][]{\lstset{style=corestyle, #1}}{}
\lstnewenvironment{svcode}  [1][]{\lstset{style=svstyle,   #1}}{}
\lstnewenvironment{cppcode} [1][]{\lstset{style=cppstyle,  #1}}{}

% Inline markup
\newcommand{\code}[1]{\texttt{\small #1}}
\newcommand{\kw}[1]{\textcolor{codekey}{\texttt{\small\bfseries #1}}}
\newcommand{\file}[1]{\textcolor{accent}{\texttt{\small #1}}}
\newcommand{\tool}[1]{\textsf{#1}}
